\section{Appendix: Weighted packets and observed-experiment converse}

The observable risk certificates are complemented by structural alternatives whose causal means are separated while their observed sample laws remain close. The construction belongs to the model class of \cref{obj:def:model-class} and compares the experiments in \cref{obj:def:observed-experiment}. Its analytic ingredient is a localized perturbation with exact cancellation under the weak-design weight.

The witness subclass uses the rank restrictions in \cref{obj:ass:rank-uniform,obj:ass:rank-treatment-independence,obj:ass:structural-threshold,obj:ass:binary-potential-outcomes}. Specifically, the auxiliary rank coordinate \(U\) is conditionally uniform on \([0,1]\) and independent of the latent dose given the stratum. A common rank generates the entire threshold schedule, and every potential outcome is binary. These restrictions give a measurable schedule with a prescribed conditional mean and place the witnesses within the bounded-outcome class. The threshold embedding and evaluation map are those of \cref{obj:def:potential-path-space}.

\subsection{Filtered Jacobi packets}

The packet construction separates localization from support. Its degree index \leanref{sym:m}{\ensuremath{m}} controls concentration and cancellation, while a fixed taper gives regularity at the boundary of the support. The filtered Jacobi sums \leanref{sym:PacketHandle}{\ensuremath{K_m,\widetilde K_m}} use an endpoint construction for positive design decay and a symmetric interior construction when the design-decay exponent is zero. The normalized packet \leanref{sym:psim}{\ensuremath{\psi_m}} is defined as follows.

\begin{definitionv}{def:packet-handle}[Filtered Jacobi packet \(\psi_m\)]
For each integer \(m\ge4\), the zero-extended filtered Jacobi packet is
\[
\psi_m(u)=
\begin{cases}
(1-u^2)^r K_m(2u^2-1)/K_m(-1),
& |u|\le1,\ \kappa>0,\\
(1-u^2)^r\widetilde K_m(u)/\widetilde K_m(0),
& |u|\le1,\ \kappa=0,\\
0,& |u|>1.
\end{cases}
\]
The fixed taper parameter and smooth spectral filter are
\[
r=4,\qquad
\eta(s)=
\begin{cases}
\exp\{-1/((s-9/8)(15/8-s))\},
& 9/8<s<15/8,\\
0,& \text{otherwise}.
\end{cases}
\]
The filter is nonnegative, supported inside \((1,2)\), and positive on \([5/4,7/4]\). Write \(\mathsf J_j^{(a,b)}\) for the standard Jacobi polynomial of degree \(j\), and define its squared norm by
\[
H_j^{a,b}
=
\int_{-1}^{1}
\bigl(\mathsf J_j^{(a,b)}(z)\bigr)^2
(1-z)^a(1+z)^b\,dz.
\]
For \(\kappa>0\), the endpoint parameter and filtered sum are
\[
b=\frac{\kappa-1}{2},\qquad
K_m(z)=\sum_{j\ge0}\eta(j/m)
\frac{\mathsf J_j^{(r,b)}(-1)\mathsf J_j^{(r,b)}(z)}
{H_j^{r,b}}.
\]
For \(\kappa=0\), the symmetric interior sum is
\[
\widetilde K_m(u)=\sum_{j\ge0}\eta(j/m)
\frac{\mathsf J_j^{(r,r)}(0)\mathsf J_j^{(r,r)}(u)}
{H_j^{r,r}}.
\]
\end{definitionv}

The two branches accommodate the parameter range of the localization estimates. For \(\kappa>0\), the endpoint parameter satisfies \(b>-1/2\). At \(\kappa=0\), the symmetric construction uses the interior point \(0\) and parameters \((r,r)\), both strictly above \(-1/2\). The fixed filter selects degrees proportional to \(m\), giving a common index for spatial concentration and weighted cancellation.

For the auxiliary polynomial identity, write \(P_j^{(a,b)}=\mathsf J_j^{(a,b)}\); both notations denote the same standard Jacobi polynomial. The parameters \(a,b\) in this identity are local polynomial parameters. For a real number \(v\) and an integer \(i\ge0\), the generalized binomial coefficient is the falling factorial divided by \(i!\), with value \(1\) at \(i=0\). The following statement records that convention together with the identity.

\begin{lemmav}{lem:jacobi-binomial}[Jacobi polynomial binomial representation]
Let \(a,b\in\mathbb R\) satisfy \(a,b>-1/2\), let \(j\ge0\) be an integer, and let \(z\in\mathbb R\). Write \(P_j^{(a,b)}=\mathsf J_j^{(a,b)}\) for the standard Jacobi polynomial. Define the generalized binomial coefficient, for \(v\in\mathbb R\) and integers \(i\ge0\), by
\[
\binom{v}{i}=\frac{\prod_{k=0}^{i-1}(v-k)}{i!},
\qquad \binom{v}{0}=1,
\]
where the empty product equals \(1\). Then
\[
P_j^{(a,b)}(z)
=2^{-j}\sum_{i=0}^{j}
\binom{j+a}{i}\binom{j+b}{j-i}
(z-1)^{j-i}(z+1)^i.
\]
\end{lemmav}

\Cref{obj:lem:jacobi-binomial} fixes the polynomial normalization used in the filtered sums. Positive squared norms and orthogonality are supplied by the standard Jacobi normalization \citep[Table 18.3.1, Jacobi row and caption]{https://dlmf.nist.gov/18.3}; the corresponding gamma-ratio asymptotic controls the degree dependence of normalization factors \citep[equation (5.11.12), positive real axis]{https://dlmf.nist.gov/5.11.E12}. The resulting packet properties are collected in the next statement.

\begin{theoremv}{thm:weighted-packet-construction}[Weighted packet construction]*
For every fixed \(\kappa\in[0,2]\), the degeneracy exponent, there exist real constants \(C>0\) and \(c>0\) such that, for every integer \(m\ge4\), the filtered-Jacobi construction has positive normalizing value:
\[
\begin{cases}
K_m(-1)>0, & \kappa>0,\\
\widetilde K_m(0)>0, & \kappa=0.
\end{cases}
\]
Here \(K_m(-1)\) and \(\widetilde K_m(0)\) are respectively the endpoint and symmetric interior normalizing values of the construction. The resulting normalized, tapered, zero-extended packet \(\psi_m\) is even and continuously differentiable on \(\mathbb R\), and satisfies
\[
\psi_m(u)=0\quad\text{for }u\notin[-1,1],
\qquad
\psi_m(0)=1,
\]
\[
|\psi_m(u)|\le C,\qquad
|\psi_m'(u)|\le Cm
\quad\text{for every }u\in\mathbb R,
\]
\[
\int_{-1}^{1}|u|^\kappa|\psi_m(u)|\,du
\le Cm^{-\kappa-1},
\qquad
\int_{-1}^{1}|u|^\kappa\psi_m(u)^2\,du
\le Cm^{-\kappa-1},
\]
and
\[
\int_{-1}^{1}|u|^\kappa u^j\psi_m(u)\,du=0
\quad\text{for every integer }0\le j<cm.
\]
\end{theoremv}
% CAUSALSMITH-CITED-SCOPE-BEGIN thm:weighted-packet-construction
\verificationfootnotetext{\textbf{Formalization scope.} This result uses the cited conclusion \citep{PetrushevXu2005} (Theorem 2.4, equation (2.14)), \citep{https://people.math.sc.edu/pencho/Publications/kpx-06-28-06-web.pdf} (Theorem 2.2, equation (2.4)), \citep{https://dlmf.nist.gov/18.3} (Table 18.3.1, Jacobi row and its caption; orthogonality and squared norm.) and \citep{https://dlmf.nist.gov/5.11.E12} (Equation 5.11.12, restricted to the positive real axis.). The cited source proof is not formalized here: Lean verifies its use and all remaining steps. If that cited conclusion is false, the portions of this result that depend on it are not certified.}
% CAUSALSMITH-CITED-SCOPE-END thm:weighted-packet-construction

The packet retains unit height at the evaluation point while its weighted absolute mass and squared mass decrease with degree. Its derivative grows at most linearly in degree, and its weighted moments vanish through an order proportional to degree. These properties serve distinct statistical purposes: normalization preserves target separation, regularity permits Hölder-admissible mean perturbations, and cancellation reduces the discrepancy after Gaussian convolution.

The localization component uses the filtered Jacobi estimate of \citet[Theorem 2.4, equation (2.14)]{PetrushevXu2005}. The local difference estimate supplies derivative control for the filtered sums \citep[Theorem 2.2, equation (2.4)]{https://people.math.sc.edu/pencho/Publications/kpx-06-28-06-web.pdf}. The taper makes the zero extension compatible with continuous differentiability. Together, these ingredients support the pointwise and weighted-integral bounds in \cref{obj:thm:weighted-packet-construction}.

Weighted cancellation is the role of Jacobi orthogonality \citep[Table 18.3.1, Jacobi row and caption]{https://dlmf.nist.gov/18.3}. In the positive-decay branch, the quadratic coordinate \(2u^2-1\) aligns the even weighted moments with the endpoint Jacobi weight; symmetry accounts for the odd moments. At zero design decay, the symmetric interior sum provides cancellation directly under the unweighted dose measure. The separate interior construction therefore supplies the same packet guarantees at that endpoint of the parameter range.

\subsection{Structural witnesses and observed distances}

We now combine the packet with independent structural coordinates. The tuning uses a localization width \(h\), a support radius \leanref{sym:ellrad}{\ensuremath{\ell}}, and a cancellation cutoff \leanref{sym:Jmom}{\ensuremath{J}}. In the inverse regimes, the relations \(h=\ell/m\) and \(J=\lfloor c_pm\rfloor\) connect these quantities to packet degree, where \(c_p>0\) is the packet's cancellation constant: its weighted moments of every order \(j<c_pm\) vanish. The grouped definition below specifies the reference and perturbed means, their structural laws, and their observed projections, and it lists the smallness and sample-size conditions under which the profiles have the asserted range and smoothness.

\begin{definitionv}{synth_5}[Structural lower-bound witness laws]
Fix \(0<\beta\le1\), \(0\le\kappa\le2\), \(0\le\sigma\le1/4\), a positive integer sample size \(n\), and the public measurable schedule space \((\mathcal S,\mathcal E)\) with threshold embedding \(\iota\). Put \(a_0=1/2\) and \(d=2\beta+\kappa+1\). For public class constants \(K=(p_{\min},c_{\mathrm{lo}},c_{\mathrm{hi}})\), require
\[
0<p_{\min}\le1/2,\qquad
c_{\mathrm{lo}}\le(\kappa+1)2^\kappa\le c_{\mathrm{hi}}.
\]
The three witness laws share independent seed coordinates \(X,T,Z,U\), where \(P\{X=0\}=P\{X=1\}=1/2\), \(Z\sim N(0,1)\), \(U\sim\operatorname{Unif}[0,1]\), and
\[
g_\kappa(t)=(\kappa+1)2^\kappa|t|^\kappa\mathbf1\{|t|\le1/2\}
\]
is the density of \(T\). Set \(A=a_0+T\).

The construction uses an amplitude \(\epsilon\ge0\), a comparison constant \(C>0\), and packet constants \(C_p,c_p>0\) such that, for every integer \(m\ge4\), the normalized filtered-Jacobi packet \(\psi_m\) of \cref{obj:def:packet-handle} is continuously differentiable on \(\mathbb R\) and satisfies
\[
|\psi_m(u)|\le C_p,\qquad
|\psi_m'(u)|\le C_pm
\quad(u\in\mathbb R),
\qquad
\int_{-1}^{1}|u|^\kappa u^j\psi_m(u)\,du=0\quad(j\in\mathbb N,\ j<c_pm).
\]
For example, the packet constants \(C_p,c_p\) of \cref{obj:thm:observed-lower} satisfy these conditions. The tuning data consist of \(h,\ell\in(0,1/4]\) and nonnegative integers \(m,J\). In the direct regime \(\sigma\le n^{-1/d}\), set
\[
h_0=n^{-1/d},\qquad h=\ell=h_0,\qquad m=J=0,
\]
and \(\delta(t)=\epsilon h_0^\beta(1-(t/h_0)^2)^2\mathbf1\{|t|\le h_0\}\), and require
\[
0<h_0\le1/4,\qquad 4\epsilon\le1,\qquad \epsilon h_0^\beta\le3/8.
\]
The bandwidth condition holds exactly when \(n\ge4^d\), so the direct-regime construction applies only to such \(n\). In either inverse regime \(\sigma>n^{-1/d}\), use the packet \(\psi_m\) and require
\[
m\ge4,\qquad h=\ell/m,\qquad J=\lfloor c_pm\rfloor\ge1,
\qquad 2\epsilon C_p\le1,\qquad \epsilon h^\beta C_p\le3/8,
\qquad \delta(t)=\epsilon h^\beta\psi_m(t/\ell).
\]
Under these requirements, the perturbation is supported on \([-\ell,\ell]\), is integrable against \(g_\kappa\), and satisfies
\[
\int_{\mathbb R}t^jg_\kappa(t)\delta(t)\,dt=0\qquad(0\le j<J);
\]
in the inverse regimes this follows from the packet cancellation because \(j<\lfloor c_pm\rfloor\) implies \(j<c_pm\). Also under these requirements, the continuous profiles
\[
\mu_\star(a)=1/2,\qquad \mu_\pm(a)=1/2\pm\delta(a-a_0)
\]
take values in \([1/8,7/8]\) and obey \(|\mu_\pm(a)-\mu_\pm(a')|\le|a-a'|^\beta\).

For \(s\in\{+,-,\star\}\), define \(P_s\) as the six-coordinate structural law of
\[
\bigl(X,A,Z,\iota(\mu_s,U),\mathbf1\{U\le\mu_s(A)\},U\bigr).
\]
Thus \(Y_s(a)=\mathbf1\{U\le\mu_s(a)\}\) simultaneously for all \(a\), \(Y_s=Y_s(A)\), and the observed projection is \((X,W,Y_s)\) with \(W=A+\sigma Z\). The retained rank \(U\) gives the common-seed representation required by the structural laws. The three laws have the same \((X,A)\) marginal and satisfy
\[
\mu_++\mu_-=2\mu_\star,\qquad
\theta(P_\pm)=1/2\pm\delta(0),\qquad \theta(P_\star)=1/2.
\]

Let \(Q_s=\mathcal L_{P_s}(X,W,Y_s)\). For every Borel set \(B\subseteq\mathbb R\) and \(x\in\{0,1\}\),
\[
Q_s\{X=x,W-a_0\in B,Y_s=1\}
=\frac12\int_{-1/2}^{1/2}g_\kappa(t)\mu_s(a_0+t)\Phi_\sigma(B-t)\,dt,
\]
with the analogous formula containing \(1-\mu_s\) when \(Y_s=0\). Here \(\Phi_\sigma=N(0,\sigma^2)\), with \(\Phi_0=\delta_0\), and \(Q_{P_s}^n=Q_s^{\otimes n}\). For \(\sigma>0\), the construction additionally requires the comparison bound below. It is a condition on the constant \(C\), not a consequence of the preceding requirements; the witnesses of \cref{obj:thm:observed-lower} satisfy it with that theorem's constant \(C\):
\[
\chi^2(Q_s,Q_\star)
\le C\sigma^{-\kappa-1}h^{2\beta+2\kappa+2}
\sum_{j=J}^{\infty}\frac{(\ell^2/\sigma^2)^j}{j!},
\qquad s\in\{+,-\}.
\]
\end{definitionv}

The calibrated design coefficient places the common density inside the public envelopes, and equal stratum probabilities satisfy the overlap restriction. Opposite perturbations of the constant reference mean create target separation while preserving the latent-design marginal. The observed formulas retain both outcome values: they compare the joint distribution of the contaminated dose and the binary outcome within each stratum.

The reference law is used to measure the discrepancy of each perturbed observed law. We make the divergence convention explicit, including its value when absolute continuity fails.

\begin{definitionv}{synth_7}[Chi-squared divergence of observed laws]
For probability measures \(Q\) and \(R\) on the observed-record space \(\{0,1\}\times\mathbb R\times[0,1]\), define
\[
\chi^2(Q,R)=
\begin{cases}
\displaystyle \int\left(\frac{dQ}{dR}-1\right)^2\,dR,
& Q\ll R,\\[4pt]
+\infty,& Q\not\ll R.
\end{cases}
\]
The integral is allowed to equal \(+\infty\). For the witness laws \(P_\pm\) and \(P_\star\), the notation \(\chi^2(P_\pm,P_\star)\) denotes
\[
\chi^2\!\left(\mathcal L_{P_\pm}(X,W,Y),
              \mathcal L_{P_\star}(X,W,Y)\right),
\qquad W=A+\sigma Z,
\]
the divergence between their single-record observed laws.
\end{definitionv}

Thus the structural-law notation in \cref{obj:synth_7} always refers to the induced observed-record comparison. The testing statement uses total variation on the observed sample space, with the following convention.

\begin{definitionv}{synth_8}[Total variation distance]
For probability measures \(\nu_1\) and \(\nu_2\) on a common measurable space \((\Omega,\mathcal F)\), define their total variation distance by
\[
\operatorname{TV}(\nu_1,\nu_2)
=\sup_{B\in\mathcal F}\bigl|\nu_1(B)-\nu_2(B)\bigr|.
\]
\end{definitionv}

Total variation controls the difference between probabilities of every observable event, including testing and interval-coverage events. The next result supplies structural witnesses with a prescribed total-variation bound and causal separation at the rate in \cref{obj:def:frontier-rate}.

\begin{theoremv}{thm:observed-lower}[Explicit observed lower witnesses]*
Fix public class constants \(K=(p_{\min},c_{\mathrm{lo}},c_{\mathrm{hi}})\) and exponents \(\beta,\kappa\) satisfying
\begin{itemize}
\item \textbf{(Class constants.)} \(0<p_{\min}\le1/2\) and \(0<c_{\mathrm{lo}}\le c_{\mathrm{hi}}\).
\item \textbf{(Exponents.)} \(0<\beta\le1\) and \(0\le\kappa\le2\).
\item \textbf{(Design calibration.)}
\[
c_{\mathrm{lo}}\le(\kappa+1)2^\kappa\le c_{\mathrm{hi}}.
\]
\end{itemize}
For every total-variation level \(\tau>0\), there exist constants \(c,C,\epsilon,C_p,c_p>0\) and an integer \(n_0\), chosen uniformly over the public potential-path space, sample size, and error scale, with the following properties.

For every integer \(m\ge4\), the packets of \cref{obj:def:packet-handle} have positive normalization:
\[
K_m(-1)>0\quad(\kappa>0),
\qquad
\widetilde K_m(0)>0\quad(\kappa=0).
\]
Each \(\psi_m\) is even and continuously differentiable on \(\mathbb R\), vanishes outside \([-1,1]\), and satisfies
\[
\psi_m(0)=1,\qquad
|\psi_m(u)|\le C_p,\qquad
|\psi_m'(u)|\le C_pm
\quad(u\in\mathbb R),
\]
\[
\int_{-1}^{1}|u|^\kappa|\psi_m(u)|\,du\le C_pm^{-\kappa-1},
\qquad
\int_{-1}^{1}|u|^\kappa\psi_m(u)^2\,du\le C_pm^{-\kappa-1},
\]
\[
\int_{-1}^{1}|u|^\kappa u^j\psi_m(u)\,du=0
\quad(j\in\mathbb N,\ j<c_pm).
\]

For every public potential-path space \((\mathcal S,\mathcal E)\) of \cref{obj:def:potential-path-space}, every integer \(n\ge n_0\), and every \(\sigma\in[0,1/4]\), there exist structural probability laws \(P_+,P_-,P_\star\) in the model class of \cref{obj:def:model-class}, with public constants \(K\) and parameters \((\beta,\kappa,\sigma)\). All three satisfy \cref{obj:ass:binary-potential-outcomes,obj:ass:rank-uniform,obj:ass:rank-treatment-independence,obj:ass:structural-threshold,obj:ass:iid}.

These laws admit the following common construction. Let \(X\) be the stratum coordinate on \(\mathcal X=\{0,1\}\), \(A\) the latent dose, \(Z\) the standardized Gaussian error, and \(U\) the rank coordinate. Write \(a_0=1/2\) for the evaluation dose and \(t=A-a_0\) for the centered latent dose. Under each law, \(X,A,Z,U\) are independent, with
\[
P(X=0)=P(X=1)=\frac12,\qquad
Z\sim N(0,1),\qquad U\sim\operatorname{Unif}[0,1].
\]
The conditional density \(g_x\) of \(A-a_0\) in either stratum is
\[
g_x(t)=
\begin{cases}
(\kappa+1)2^\kappa|t|^\kappa,& |t|\le1/2,\\
0,& |t|>1/2,
\end{cases}
\qquad x\in\mathcal X.
\]
There are continuous profiles \(\mu_+,\mu_-:[0,1]\to[0,1]\) and a reference profile \(\mu_\star(a)=1/2\). Under \(P_+\), \(P_-\), and \(P_\star\), respectively, the random schedule and realized outcome are
\[
Y(\cdot)=\iota(\mu_+,U),\quad
Y(\cdot)=\iota(\mu_-,U),\quad
Y(\cdot)=\iota(\mu_\star,U),
\qquad
Y=e(Y(\cdot),A).
\]
Here \(e\) and \(\iota\) are the evaluation and threshold-schedule maps of \cref{obj:def:potential-path-space}. In particular, the latent-design marginals agree:
\[
\mathcal L_{P_+}(X,A)=
\mathcal L_{P_-}(X,A)=
\mathcal L_{P_\star}(X,A).
\]

Writing
\[
\theta(P)=\mathbb E_P[Y(a_0)]
\]
for the population causal mean, and using the frontier rate \(\rho_{n,\sigma}\) of \cref{obj:def:frontier-rate} and the observed product experiment \(Q_P^n\) of \cref{obj:def:observed-experiment}, the witnesses satisfy
\[
|\theta(P_+)-\theta(P_-)|\ge c\rho_{n,\sigma},
\qquad
\operatorname{TV}(Q_{P_+}^n,Q_{P_-}^n)\le\tau.
\]

The same profiles have common tuning data \(h,\ell\in\mathbb R\) and \(m,J\in\mathbb N\), where \(h\) is the localization width, \(\ell\) the support radius, \(m\) the packet degree, and \(J\) the cancellation cutoff. Define the direct-regime bandwidth
\[
h_0=n^{-1/(2\beta+\kappa+1)}.
\]
The tuning satisfies
\[
0<h\le1/4,\qquad 0<\ell\le1/4,
\]
\[
\sigma\le h_0\ \Longrightarrow\
h=h_0,\quad \ell=h,\quad J=0,
\]
\[
h_0<\sigma\ \Longrightarrow\
m\ge4,\quad h=\ell/m,\quad
J=\lfloor c_pm\rfloor,\quad J\ge1.
\]
Define the mean perturbation \(\delta\) by
\[
\delta(t)=
\begin{cases}
\epsilon h_0^\beta\bigl(1-(t/h_0)^2\bigr)^2,
  &\sigma\le h_0,\ |t|\le h_0,\\
0,&\sigma\le h_0,\ |t|>h_0,\\
\epsilon h^\beta\psi_m(t/\ell),&h_0<\sigma.
\end{cases}
\]
For every \(a\in[0,1]\),
\[
\mu_\pm(a)=\frac12\pm\delta(a-a_0).
\]
Moreover,
\[
\delta(t)=0\quad(|t|>\ell),
\qquad
\int_{\mathbb R}t^jg_x(t)\delta(t)\,dt=0
\quad(x\in\mathcal X,\ j\in\mathbb N,\ j<J).
\]

For \(\sigma>0\), let \(\chi^2(P_\pm,P_\star)\) denote the chi-squared divergence between the respective one-observation laws induced by \(P_\pm\) and \(P_\star\) in \cref{obj:def:observed-experiment}. Both signs satisfy
\[
\chi^2(P_\pm,P_\star)
\le
C\sigma^{-\kappa-1}h^{2\beta+2\kappa+2}
\sum_{j=J}^{\infty}\frac{(\ell^2/\sigma^2)^j}{j!}.
\]
\end{theoremv}
% CAUSALSMITH-CITED-SCOPE-BEGIN thm:observed-lower
\verificationfootnotetext{\textbf{Formalization scope.} This result uses the cited conclusion \citep{PetrushevXu2005} (Theorem 2.4, equation (2.14)), \citep{https://people.math.sc.edu/pencho/Publications/kpx-06-28-06-web.pdf} (Theorem 2.2, equation (2.4)), \citep{https://dlmf.nist.gov/18.3} (Table 18.3.1, Jacobi row and its caption; orthogonality and squared norm.) and \citep{https://dlmf.nist.gov/5.11.E12} (Equation 5.11.12, restricted to the positive real axis.). The cited source proof is not formalized here: Lean verifies its use and all remaining steps. If that cited conclusion is false, the portions of this result that depend on it are not certified.}
% CAUSALSMITH-CITED-SCOPE-END thm:observed-lower

The witnesses change the causal response while keeping the distribution of stratum and latent dose fixed. Their opposite mean perturbations preserve a common reference profile, and the threshold representation carries those profiles into legal potential-outcome schedules. The comparison concerns the complete observed record, including its outcome mark. Consequently, the separation and total-variation clauses address the same experiment used by the estimation and interval criteria.

The analytic inputs to the packet component of \cref{obj:thm:observed-lower} are the localization estimate \citep[Theorem 2.4, equation (2.14)]{PetrushevXu2005}, the local difference estimate \citep[Theorem 2.2, equation (2.4)]{https://people.math.sc.edu/pencho/Publications/kpx-06-28-06-web.pdf}, Jacobi orthogonality and squared norms \citep[Table 18.3.1, Jacobi row and caption]{https://dlmf.nist.gov/18.3}, and the gamma-ratio asymptotic \citep[equation (5.11.12), positive real axis]{https://dlmf.nist.gov/5.11.E12}. Their roles remain normalization, localization, derivative control, and weighted cancellation.

\subsection{Regime-specific observed-law comparison}

The direct and inverse constructions in \cref{obj:thm:observed-lower} balance target separation against observable discrimination. In the direct regime, \(\sigma\le h_0\), the compact bump has support radius \(h_0=n^{-1/d}\) and height \(\epsilon h_0^\beta\). The witness separation is therefore \(2\epsilon h_0^\beta\), yielding the direct rate \(n^{-\beta/d}\). The theorem's observed-sample total-variation bound covers this entire regime, including \(\sigma=0\).

For \(h_0<\sigma\), packet normalization gives separation \(2\epsilon h^\beta\). The support radius and localization width differ by the packet degree, and cancellation removes all weighted moments below \(J\). In the displayed chi-squared comparison, the surviving series begins at that cutoff; its argument \(\ell^2/\sigma^2\) records the relation between perturbation support and measurement error. The comparison therefore incorporates both the cancellation order and the permitted support radius.

Across the intermediate range
\[
n^{-1/d}<\sigma\le L_n^{-1/2},
\]
the separation guarantee has order
\[
\left\{\frac{\sigma}{\sqrt{\log(e+n\sigma^d)}}\right\}^{\beta}.
\]
Across the compact-support range
\[
L_n^{-1/2}<\sigma\le1/4,
\]
it has order
\[
\left\{\frac{\log(e+\sigma^2L_n)}{L_n}\right\}^{\beta}.
\]
These are the two inverse branches of \cref{obj:def:frontier-rate}. Both are realized by packets supported within the stated radius bound, with a cancellation cutoff proportional to degree. The common total-variation guarantee in \cref{obj:thm:observed-lower} completes the observed-law comparison uniformly over the error-scale range and the declared potential-path spaces, beyond its common sample-size threshold.

\subsection{Risk, honest length, and transition consequences}

Let \(\Delta=|\theta(P_+)-\theta(P_-)|\) denote the witness target separation. The two-point testing bound converts a separated pair of targets and close sample laws into an absolute-error lower bound \citep[Section 2.2, Theorem 2.2]{Tsybakov2009}. With any fixed \(\tau\in(0,1)\), the witnesses in \cref{obj:thm:observed-lower} give
\[
R_n(\sigma)\ge \frac{\Delta}{4}(1-\tau)
\ge c_R\rho_{n,\sigma},
\]
for a positive constant \(c_R\), every sufficiently large sample size, every \(\sigma\in[0,1/4]\), and every public potential-path space in \cref{obj:def:potential-path-space}. The fixed parameters satisfy the design calibration and exponent restrictions of \cref{obj:thm:observed-lower}. This lower bound applies to the possibly randomized estimators allowed by \cref{obj:def:minimax-risk}.

For honest connected intervals, fix \(\alpha\in(0,1/2)\) and choose \(\tau\in(0,1-2\alpha)\). Uniform coverage at the two witness laws and their total-variation bound give the corresponding expected-length consequence
\[
H_n(\sigma)\ge \Delta(1-2\alpha-\tau)
\ge c_H\rho_{n,\sigma},
\]
with \(c_H>0\), under the same uniform large-sample quantifiers. Connectedness makes an interval containing both witness targets have length at least their separation. The factor \(1-2\alpha-\tau\) records the coverage requirement and the observable discrepancy of the pair. These consequences supply the lower bounds in \cref{obj:thm:uniform-frontier}; the observable upper bounds are provided by \cref{obj:thm:observable-upper,obj:thm:finite-honesty}.

The transition comparisons in \cref{obj:thm:uniform-frontier} connect the three descriptions throughout fixed multiplicative neighborhoods of their thresholds. Using the transition thresholds \(a_n,b_n\) and comparison scales \(D_n,F_n(\sigma),P_n(\sigma)\) defined there, for every fixed multiplicative factor \(M\ge1\),
\[
F_n(\sigma)\asymp D_n
\qquad
\bigl(\sigma\in[a_n/M,Ma_n]\bigr),
\]
and
\[
F_n(\sigma)\asymp P_n(\sigma)
\qquad
\bigl(\sigma\in[b_n/M,Mb_n]\bigr),
\]
once the sample size exceeds the stated window threshold. The constants and threshold may depend on \(M\) and the fixed parameters. Raising these comparisons to the smoothness exponent shows that neighboring rate formulas give comparable risk and honest length throughout each window.

At the error-free endpoint, \cref{obj:prop:error-free-reduction} gives
\[
R_n(0)\asymp H_n(0)\asymp n^{-\beta/(2\beta+\kappa+1)}.
\]
For every fixed \(\sigma\in(0,1/4]\), the final conclusion of \cref{obj:thm:uniform-frontier} instead gives
\[
R_n(\sigma)\asymp H_n(\sigma)
\asymp
\left(\frac{\log\log n}{\log n}\right)^\beta.
\]
The first comparison uses constants determined by the calibrated class and fixed coverage level; the second permits dependence on the fixed positive error scale. Both retain uniformity over the declared public potential-path spaces. The observed-law witnesses thus support the direct endpoint, both inverse regimes, and the transition comparisons within one structural model.

\paragraph{Verification note.}
The theorem statements and proofs are machine-checked in Lean 4. The definition blocks written for exposition (\cref{obj:synth_1,obj:synth_2,obj:synth_3,obj:synth_4,obj:synth_5,obj:synth_6,obj:synth_7,obj:synth_8}) restate in prose the Lean definitions and constructions that the theorems use; these prose blocks are not themselves machine-checked statements. The model assumptions specify the structural class under which the conclusions hold. The theorem-local verification notes identify the cited analytic conclusions, whose published source proofs remain external, and record the scope of their use.
